% Answers to Chapter 6, from lectures/L06/appendix/c_solutions.md; the self-test from
% lectures/L06/README.md. Then the answers to Övningsdugga 1, from
% lectures/L06/appendix/practice_test/solutions.md.

\facitavsnitt{c6}{Kapitel 6}{Andragradsekvationer}
\begin{facit}

\svar{1.1} \sv{a} $10$ \sv{b} $4 \cdot 10^6$ \sv{c} $2 \times 10^3 = 2\,000$

\svar{1.2} \sv{a} $I \approx 7{,}02 \times 10^{-1}\,\text{mA} = 0{,}702\,\text{mA}$
\sv{b} $P \approx 2{,}32\,\text{mW}$

\svar{1.3} \sv{a} $12$ \sv{b} $0{,}07$ \sv{c} $20\,\text{V}$

\svar{2.1} \sv{a} $I = \pm 0{,}4\,\text{A}$
\sv{b} Den negativa roten anger bara att strömmen går åt andra hållet; effekten $RI^2$ är
densamma, så strömmens storlek är $I = 0{,}4\,\text{A}$.

\svar{2.2} \sv{a} $x_1 = 3$, $x_2 = 2$ \sv{b} $x_1 = 2$, $x_2 = -5$
\sv{c} $x_1 = x_2 = 3$ (dubbelrot)

\svar{2.3} \sv{a} $x_1 = 3$, $x_2 = 0{,}5$ \sv{b} $x_1 = \frac{1}{3}$, $x_2 = -2$

\svar{2.4} \sv{a} $R_1^2 - 13R_1 + 36 = 0$ \sv{b} $9\,\Omega$ och $4\,\Omega$

\svar{3.1} \sv{a} $x = \pm 7$ \sv{b} $x = \pm 3$ \sv{c} $x = \pm 5$
\sv{d} Inga reella lösningar, eftersom $x^2 = -4$. \sv{e} $x = \pm 6$

\svar{3.2} \sv{a} $x_1 = 0$, $x_2 = 7$ \sv{b} $x = \pm 4$ \sv{c} $x_1 = -2$, $x_2 = -3$
\sv{d} $x_1 = 0$, $x_2 = 4$ \sv{e} $x_1 = x_2 = 2$ (dubbelrot)

\svar{3.3} \sv{a} $x_1 = 2$, $x_2 = -4$ \sv{b} $x_1 = 6$, $x_2 = 2$ \sv{c} $x_1 = 3$, $x_2 = -4$
\sv{d} $x_1 = 5$, $x_2 = -3$

\svar{3.4} \sv{a} $x_1 = 3$, $x_2 = 2$ \sv{b} $x_1 = 1$, $x_2 = -5$ \sv{c} $x_1 = 3$, $x_2 = 1$

\svar{3.5} \sv{a} $D = 16$, två reella rötter \sv{b} $D = 0$, en dubbelrot
\sv{c} $D = -11$, inga reella rötter \sv{d} $D = 0$, en dubbelrot
\sv{e} $D = 16$, två reella rötter

\svar{3.6} \sv{a} $x_1 = 1$, $x_2 = -5$ \sv{b} $x_1 = 7$, $x_2 = 3$
\sv{c} $x = -3 \pm \sqrt{5}$, alltså $x_1 \approx -0{,}76$ och $x_2 \approx -5{,}24$

\svar{3.7} \sv{a} $p = -2$, $q = -8$: $x^2 - 2x - 8 = 0$
\sv{b} $x = 1 \pm 3$, alltså $4$ och $-2$.
\sv{c} $4 + 5 = 9 = -p$ och $4 \cdot 5 = 20 = q$, så rötterna stämmer.
\sv{d} Ja: $x^2 - 6x + 9 = 0$, med $D = 0$.

\svar{3.8} \sv{a} $0{,}5\,\text{A}$ \sv{b} $30\,\text{mA}$ \sv{c} Cirka $33\,\text{mA}$
\sv{d} Den negativa roten betyder bara motsatt strömriktning; effekten blir densamma eftersom $I$
kvadreras, och det är strömmens storlek som efterfrågas.

\svar{3.9} \sv{a} $10\,\text{V}$ \sv{b} Cirka $15{,}8\,\text{V}$ \sv{c} $52{,}9\,\Omega$

\svar{3.10} \sv{a} $12\,\Omega$ och $8\,\Omega$ \sv{b} $8\,\Omega$ och $7\,\Omega$
\sv{c} Nej. $x^2 - 10x + 30 = 0$ har $D = -20 < 0$ och alltså inga reella rötter.

\svar{3.11} \sv{a} Produkt genom summa ger $R_1 R_2 = (R_1 \parallell R_2)(R_1 + R_2)
= 2{,}4 \times 10 = 24\,\Omega^2$.
\sv{b} $x^2 - 10x + 24 = 0$ ger $6\,\Omega$ och $4\,\Omega$.
\sv{c} $6 + 4 = 10\,\Omega$ och $\frac{6 \times 4}{6 + 4} = 2{,}4\,\Omega$.

\svar{3.12} \sv{a} $P = R\left(\frac{12}{4 + R}\right)^2 = 8$ ger $18R = (4 + R)^2$, alltså
$R^2 - 10R + 16 = 0$.
\sv{b} $R = 2\,\Omega$ eller $R = 8\,\Omega$
\sv{c} $R = 2\,\Omega$: $I = 2\,\text{A}$ och $P = 8\,\text{W}$. $R = 8\,\Omega$: $I = 1\,\text{A}$
och $P = 8\,\text{W}$.
\sv{d} $I = 1{,}5\,\text{A}$ och $P = 9\,\text{W}$, mer än de två lösningarnas $8\,\text{W}$:
effekten i $R$ är störst när $R = R_1$ (effektanpassning).

\svar{3.13} \sv{a} Falskt: det beror på diskriminanten, som kan ge två, en eller ingen reell rot.
\sv{b} Falskt: även $x = -4$ är en lösning.
\sv{c} Falskt: ekvationen måste först normeras, $x^2 + 2x - 3 = 0$.
\sv{d} Sant: rottermen blir noll, och båda rötterna blir $-\frac{p}{2}$.
\sv{e} Sant: det är nollproduktmetoden.
\sv{f} Sant: $x^2 = -4$, och ingen reell kvadrat är negativ.

\svar[Testa dig själv]{T} \sv{1} $2x(x - 4) = 0$, alltså $x_1 = 0$ och $x_2 = 4$.
\sv{2} $I = 0{,}5$ (ampere). Den negativa roten förkastas: den anger bara motsatt strömriktning,
och effekten $RI^2$ blir densamma.

\end{facit}

\facitavsnitt{od1}{Övningsdugga 1}{Aritmetik, algebra, potenser och ekvationer}
\begin{facit}

\svar{1} $R_2 \approx 8\,\text{k}\Omega$ (oavrundat $8{,}18\,\text{k}\Omega$)

\svar{2} $c(c^2 + 1) = c^3 + c$

\svar{3} $x = -\frac{17}{7}$, $y = -\frac{19}{7}$

\svar{4} $k = 0{,}5$, $m = -1{,}5$

\svar{5} $P \approx 2{,}9 \cdot 10^{-2}\,\text{W}$ (cirka $29\,\text{mW}$)

\svar{6} $x_1 \approx 0{,}21$, $x_2 \approx -27{,}54$

\end{facit}
